\subsection{Evaluasi terhadap Rumusan Masalah}
\label{subsec:evaluasi-rumusan-masalah}

Evaluasi terhadap rumusan masalah dilakukan dengan melihat apakah implementasi dapat menjawab masalah utama yang telah diidentifikasi, yaitu keterbatasan kontrol akses granular, kebutuhan segmentasi data RME, dan kebutuhan delegasi akses yang tetap membatasi cakupan hak pemegang token.

Pertama, implementasi menunjukkan bahwa Macaroons dapat digunakan sebagai token otorisasi yang lebih ekspresif dibandingkan token akses yang hanya bergantung pada peran. \textit{Caveats} memungkinkan hak akses dinyatakan dalam bentuk batasan yang lebih rinci, seperti kategori data, aksi, masa berlaku, dan konteks delegasi. Dengan demikian, sistem dapat membedakan kebutuhan akses dokter, laboratorium, dan apoteker secara lebih spesifik.

Kedua, segmentasi data RME memungkinkan sistem menegakkan prinsip \textit{need-to-know}. Karena setiap segmen memiliki metadata kategori, \textit{Server} PRE dapat menolak permintaan terhadap segmen yang tidak tercakup oleh token. Pendekatan ini mengurangi risiko aktor memperoleh seluruh data RME ketika sebenarnya hanya membutuhkan bagian tertentu untuk menjalankan tugas klinisnya.

Ketiga, atenuasi \textit{offline} pada Macaroons memungkinkan delegasi akses dilakukan tanpa penerbitan ulang token oleh pasien. Hal ini menjawab kebutuhan alur layanan klinis yang dinamis, misalnya ketika dokter perlu mendelegasikan sebagian akses kepada laboratorium atau apoteker. Pada saat yang sama, sifat \textit{append-only} Macaroons memastikan token turunan tidak dapat menghapus batasan lama sehingga cakupan akses tetap sama atau lebih sempit. Implementasi juga menunjukkan bahwa token turunan dapat dibuat di \textit{Hospital Client} tanpa kunci akar, lalu digunakan oleh penerima delegasi melalui \textit{Server} PRE.

Berdasarkan evaluasi tersebut, rancangan yang diimplementasikan dapat menjawab rumusan masalah pada ruang lingkup prototipe. Namun, hasil ini belum dapat dianggap sebagai validasi produksi karena implementasi belum diuji pada lingkungan rumah sakit nyata, volume data besar, integrasi lintas institusi, maupun uji keamanan formal end-to-end. Selain itu, beberapa bagian enforcement masih parsial, terutama verifikasi \texttt{delegation\_signature}, endpoint legacy, dan validasi kesegaran waktu pada bukti \textit{wallet}.
